% 2.2 — brief en documento/investigacion/2-2-aprendizaje-automatico.md
\subsection{APRENDIZAJE AUTOMÁTICO}
\subsubsection{TIPOS DE APRENDIZAJE}

El aprendizaje automático (\textit{machine learning}) es la rama de la inteligencia artificial que estudia los algoritmos capaces de mejorar su desempeño de manera automática a partir de la experiencia. \textcite{Mitchell1997} lo formaliza con tres elementos: se dice que un programa aprende de una experiencia $E$ respecto de una clase de tareas $T$ y una medida de desempeño $P$ si su desempeño en $T$, medido por $P$, mejora con $E$. Trasladada al presente proyecto, esa formalización ubica la experiencia en el histórico operativo de más de diez años de ISI Mustang y la tarea en anticipar desviaciones en sus proyectos, mientras que la medida de desempeño está dada por las métricas que se describen en la subsección siguiente.

La forma en que esa experiencia se presenta al algoritmo define los paradigmas clásicos del campo. \textcite{Geron2022} organiza los sistemas de aprendizaje automático según la supervisión que reciben durante el entrenamiento en tres familias: supervisado, no supervisado y por refuerzo. En el aprendizaje supervisado cada observación de entrenamiento viene acompañada de la respuesta correcta, la etiqueta o valor objetivo, y el algoritmo aprende una función que asocia las entradas con esa salida para predecirla sobre observaciones nuevas \parencite{Hastie2009}. En el aprendizaje no supervisado no existe tal respuesta: el algoritmo recibe únicamente las entradas y busca estructura en ellas, como grupos de observaciones similares o direcciones de mayor variabilidad \parencite{James2021}. El aprendizaje por refuerzo, por último, no parte de un conjunto de datos fijo sino de un agente que interactúa con un entorno y ajusta su comportamiento según las recompensas que recibe \parencite{Mitchell1997}.

Dentro del aprendizaje supervisado, la naturaleza de la variable de salida distingue dos tipos de problema. Cuando la salida es cuantitativa (un monto, una duración, un porcentaje) se habla de regresión; cuando es cualitativa o categórica (una clase entre un conjunto finito) se habla de clasificación \parencite{Hastie2009, James2021}. La distinción no es solo nominal: determina qué algoritmos son aplicables, qué función de pérdida se optimiza y con qué métricas se evalúa el resultado.

Los tres casos de uso del módulo predictivo se plantean en el paradigma supervisado: la predicción de desviación en certificaciones como un problema de clasificación, y la estimación del sobrecosto y del retraso en la fecha de cierre como problemas de regresión. El planteo supone que el histórico de la empresa contiene, para cada proyecto cerrado, el resultado que se quiere predecir: cuánto costó respecto de lo presupuestado, cuándo se cerró respecto de la fecha comprometida y cómo evolucionaron sus certificaciones frente a lo planificado. La existencia, completitud y recuperabilidad de esas respuestas para una proporción suficiente de proyectos se examina en el análisis de datos del Capítulo III, del que dependen también los detalles de las variables y del protocolo de evaluación que esta sección presenta en términos generales. El aprendizaje no supervisado y el aprendizaje por refuerzo quedan fuera del alcance del proyecto.

\subsubsection{APRENDIZAJE SUPERVISADO}

Formalmente, el aprendizaje supervisado parte de un conjunto de entrenamiento compuesto por pares $(x_i, y_i)$, donde $x_i$ es un vector de variables predictoras y $y_i$ la respuesta observada, y busca una función $f$ tal que $f(x)$ aproxime a $y$ con el menor error posible sobre observaciones que no participaron del entrenamiento \parencite{Hastie2009}. Esta última condición es la que separa aprender de memorizar. \textcite{James2021} distinguen el error de entrenamiento, que mide el ajuste sobre los datos vistos, del error de prueba, que mide la capacidad de generalizar; un modelo puede reducir el primero indefinidamente aumentando su complejidad, pero a partir de cierto punto el segundo empieza a crecer. Ese fenómeno es el sobreajuste, y se explica por el compromiso entre sesgo y varianza: los modelos simples tienen sesgo alto porque no capturan la estructura de los datos, mientras que los modelos muy flexibles tienen varianza alta porque se ajustan al ruido de la muestra particular con la que fueron entrenados \parencite{Hastie2009, James2021}.

Como el error de prueba no puede calcularse sobre los datos de entrenamiento, se requiere un procedimiento para estimarlo. La práctica básica consiste en separar una parte de los datos como conjunto de prueba, que se reserva y solo se usa al final para evaluar el modelo elegido \parencite{Geron2022}. Para seleccionar entre modelos o ajustar hiperparámetros sin contaminar ese conjunto se emplea la validación cruzada de $k$ particiones (\textit{k-fold}): los datos de entrenamiento se dividen en $k$ bloques, el modelo se entrena $k$ veces dejando fuera un bloque distinto en cada corrida y el error se promedia sobre las $k$ evaluaciones \parencite{James2021}. En un estudio empírico de referencia sobre selección de clasificadores, \textcite{Kohavi1995} concluye que la validación cruzada estratificada de diez particiones es el mejor método para selección de modelos, incluso cuando la capacidad de cómputo permite usar más particiones, y advierte que la variante \textit{leave-one-out}, aunque casi insesgada, tiene alta varianza y produce estimaciones poco confiables. La estratificación conserva en cada bloque la proporción de clases del conjunto completo; es propia de los problemas de clasificación e importa cuando una de las clases es minoritaria.

La evaluación del modelo depende del tipo de problema. En clasificación, las medidas usuales derivan de la matriz de confusión, que cruza la clase real con la clase predicha \parencite{Sokolova2009, Geron2022}; de ella se obtienen la exactitud, la precisión, la sensibilidad o \textit{recall} y la medida F1, definidas en el cuadro~\ref{tab:metricas-supervisado}. \textcite{Sokolova2009} analizan qué cambios en la matriz de confusión dejan invariante cada medida y muestran que estas cuatro reaccionan de forma distinta ante cambios en la distribución de clases, por lo que conviene reportarlas juntas en lugar de limitarse a la exactitud.

En regresión, el punto de partida es el error cuadrático medio (MSE), que al elevar al cuadrado las diferencias entre valor observado y predicho da más peso a los errores grandes \parencite{James2021}; su raíz, el RMSE, se expresa en las unidades de la variable, mientras que el error absoluto medio (MAE), al promediar diferencias en valor absoluto, penaliza menos esos errores. \textcite{Chai2014} sostienen que el RMSE es más apropiado cuando se espera que los errores sigan una distribución gaussiana y que evaluar un modelo requiere una combinación de métricas y no una sola. El coeficiente de determinación $R^2$ mide la proporción de la variabilidad de la respuesta que el modelo explica; toma la forma de una proporción y no depende de la escala de la variable \parencite{James2021}, aunque qué valor se considera bueno depende de la aplicación. El cuadro~\ref{tab:metricas-supervisado} resume las métricas y su uso en el proyecto.

\begin{table}[H]
  \centering
  \caption{Métricas de evaluación de modelos supervisados y su uso en el proyecto}\label{tab:metricas-supervisado}
  \begin{tblr}{colspec={X[0.57,l] X[1.05,l] X[1.37,l]}}
    \textbf{Métrica} & \textbf{Qué mide} & \textbf{Uso en el proyecto} \\
    \SetCell[c=3]{l} \textit{Clasificación} & & \\
    Matriz de confusión & Cruce de clase real y clase predicha (VP, VN, FP, FN) & Base de las demás métricas; muestra qué tipo de error comete el clasificador de desviación en certificaciones \\
    Exactitud & $(\mathrm{VP}+\mathrm{VN})/\mathrm{total}$ & Referencia general; insuficiente si las clases están desbalanceadas \\
    Precisión & $\mathrm{VP}/(\mathrm{VP}+\mathrm{FP})$ & Proporción de alertas de desviación que son correctas \\
    Sensibilidad & $\mathrm{VP}/(\mathrm{VP}+\mathrm{FN})$ & Proporción de desviaciones reales que el modelo detecta \\
    F1 & Media armónica de precisión y sensibilidad & Resumen único cuando ambas importan \\
    \SetCell[c=3]{l} \textit{Regresión} & & \\
    RMSE & Raíz del promedio de errores al cuadrado & Error en unidades de la variable; penaliza sobrecostos o retrasos grandes mal estimados \\
    MAE & Promedio de errores en valor absoluto & Error típico esperado, menos sensible a valores extremos \\
    $R^2$ & Proporción de la variabilidad explicada por el modelo & Ajuste de cada regresor sobre su propia respuesta, independiente de la escala; complementa al RMSE y al MAE \\
  \end{tblr}
  \fuente{Elaboración propia, 2026, en base a \textcite{Sokolova2009}, \textcite{James2021} y \textcite{Chai2014}.}
\end{table}

Con estas definiciones, el protocolo de evaluación del módulo predictivo es el siguiente: el clasificador de desviación en certificaciones se evalúa con exactitud, precisión, sensibilidad, F1 y matriz de confusión; los regresores de sobrecosto y de retraso en la fecha de cierre, con RMSE, MAE y $R^2$. En los tres casos la selección de hiperparámetros se realiza con validación cruzada de $k$ particiones sobre el conjunto de entrenamiento (estratificada solo en el clasificador) y se reserva un conjunto de prueba independiente para la evaluación final.

\subsubsection[ALGORITMOS: GRADIENT BOOSTING Y XGBOOST]{ALGORITMOS DE APRENDIZAJE SUPERVISADO: GRADIENT BOOSTING Y XGBOOST}

El modelo base de los algoritmos empleados en este proyecto es el árbol de decisión. Un árbol particiona recursivamente el espacio de las variables predictoras mediante reglas binarias del tipo «presupuesto mayor que cierto umbral» y asigna a cada región resultante una predicción constante: la clase mayoritaria en clasificación o el promedio de las respuestas en regresión \parencite{Hastie2009, James2021}. El procedimiento CART de \textcite{Breiman1984} es la formulación de referencia de este tipo de árboles binarios. Los árboles son fáciles de interpretar y, en su formulación teórica, pueden dividir directamente sobre variables categóricas; las implementaciones actuales de \textit{gradient boosting} admiten esa capacidad de forma nativa como opción, XGBoost desde su versión 1.5 y Scikit-learn en sus estimadores basados en histogramas. En la configuración adoptada en este proyecto, sin embargo, los predictores categóricos se codifican de forma explícita antes del entrenamiento, como se describe en la subsección siguiente, lo que deja la representación de las variables bajo el control del pipeline de datos y no de la biblioteca. Por otra parte, un árbol aislado tiene alta varianza (pequeños cambios en los datos producen árboles muy distintos) y su capacidad predictiva suele ser inferior a la de otros métodos \parencite{James2021}.

Los métodos de ensamble corrigen esa debilidad combinando muchos árboles, con dos estrategias distintas. En el \textit{bagging}, los árboles se entrenan en paralelo sobre remuestreos de los datos y sus predicciones se promedian, lo que reduce la varianza; los bosques aleatorios de \textcite{Breiman2001} parten de esa idea y añaden una selección aleatoria de variables en cada división. En el \textit{boosting}, en cambio, los árboles se construyen de forma secuencial: cada nuevo árbol se ajusta a lo que los anteriores todavía no explican, de modo que el ensamble mejora por etapas \parencite{Hastie2009, James2021}.

\textcite{Friedman2001} da a esta idea su formulación general al plantear la estimación de la función como un problema de optimización numérica en el espacio de funciones: en cada etapa se ajusta un árbol de regresión al gradiente negativo de la función de pérdida evaluada sobre el modelo acumulado, lo que equivale a un descenso por gradiente en el que cada paso es un árbol. El marco admite distintas funciones de pérdida (mínimos cuadrados, desviación absoluta o Huber para regresión, y verosimilitud logística para clasificación), de modo que el mismo algoritmo sirve para los problemas de regresión y de clasificación del proyecto. Sus hiperparámetros principales son el número de árboles, la tasa de aprendizaje o \textit{shrinkage}, que reduce la contribución de cada árbol, y la profundidad de los árboles, que controla el orden de las interacciones entre variables \parencite{James2021}; el submuestreo de observaciones en cada etapa agrega una regularización adicional \parencite{Hastie2009}.

XGBoost es la implementación de \textit{gradient boosting} de árboles que se prevé utilizar. \textcite{Chen2016} la presentan como un sistema escalable, ampliamente usado para obtener resultados de vanguardia en competencias de aprendizaje automático, y describen sus aportes: un algoritmo \textit{sparsity-aware} para datos dispersos, un esquema de cuantiles ponderados para construir árboles de forma aproximada, y optimizaciones de caché, compresión y particionado que le permiten escalar a miles de millones de ejemplos con menos recursos que otros sistemas.

Dos características de XGBoost interesan al diseño del módulo predictivo \parencite{Chen2016}. La primera es el tratamiento de valores ausentes: el algoritmo de búsqueda de divisiones \textit{sparsity-aware} incorpora en cada nodo una dirección por defecto, aprendida de los datos, hacia la que se envían las observaciones que carecen del valor de la variable de división, lo que permite entrenar sin imputar previamente los valores ausentes. La segunda es la función objetivo regularizada, que suma a la pérdida un término de complejidad que penaliza el número de hojas y la magnitud de los pesos de cada árbol; sus coeficientes son hiperparámetros adicionales a los del \textit{boosting} ya mencionados y se ajustan mediante validación cruzada.

La elección de \textit{gradient boosting} frente a alternativas más recientes se apoya en evidencia sobre datos tabulares. \textcite{Grinsztajn2022} comparan modelos de árboles y redes neuronales profundas sobre 45 conjuntos de datos tabulares y concluyen que los primeros siguen siendo el estado del arte en conjuntos de tamaño medio, del orden de diez mil observaciones, aun sin considerar su ventaja en velocidad; atribuyen el resultado a sesgos inductivos de los árboles, como la robustez frente a variables no informativas y la capacidad de aprender funciones irregulares, que las redes no comparten. \textcite{ShwartzZiv2022} llegan a la misma conclusión al enfrentar XGBoost con modelos profundos diseñados específicamente para datos tabulares: XGBoost los supera incluso en los conjuntos de datos usados por los artículos que los proponen y requiere mucho menos ajuste. Ambos resultados se obtienen sobre datos tabulares, la forma que toma el histórico de la empresa, y respaldan la elección de XGBoost para los tres casos de uso. En la implementación se emplea su interfaz compatible con Scikit-learn, la biblioteca de aprendizaje automático en Python presentada por \textcite{Pedregosa2011}. Cabe una salvedad: los conjuntos evaluados en esos estudios son del orden de miles de observaciones, mientras que el histórico de la empresa, con un vector de variables por proyecto, probablemente reúne bastantes menos, lo que condiciona hasta qué punto esa evidencia es trasladable.

\subsubsection{FEATURE ENGINEERING}

La ingeniería de características (\textit{feature engineering}) es el conjunto de operaciones que transforman los datos crudos en las variables predictoras que recibe el modelo. \textcite{Kuhn2019} la sitúan dentro del proceso de modelado predictivo, junto con la división de datos, el remuestreo, el ajuste de hiperparámetros y la evaluación, y sostienen que la representación de los predictores condiciona el desempeño tanto como la elección del algoritmo. \textcite{Geron2022} muestra el mismo proceso desde la práctica con Scikit-learn: limpieza, codificación de variables categóricas, escalado y encadenamiento de esas transformaciones en \textit{pipelines} que se ajustan únicamente con datos de entrenamiento, para evitar que información del conjunto de prueba se filtre al modelo.

Tres grupos de técnicas son relevantes para este proyecto. El primero es la codificación de predictores categóricos, como el cliente, el tipo de proyecto o el rubro, que en este proyecto se realiza de forma explícita en el pipeline según la decisión indicada en la subsección anterior: \textcite{Kuhn2019} describen la codificación \textit{one-hot} o por variables indicadoras, el agrupamiento de categorías poco frecuentes y las codificaciones supervisadas, que reemplazan cada categoría por un estadístico de la respuesta, útiles cuando la variable tiene muchos niveles. El segundo es la transformación de predictores numéricos, que incluye el centrado y escalado, las transformaciones para corregir asimetría y la discretización en intervalos. El tercero, y el más sensible en un histórico operativo, es el manejo de datos faltantes: los mismos autores recomiendan visualizar primero el patrón de ausencia y luego decidir entre eliminar observaciones o variables, imputar los valores o emplear modelos tolerantes a faltantes. Esta última opción se corresponde con el tratamiento de valores ausentes de XGBoost descrito en la subsección anterior.

La selección de variables completa el proceso. \textcite{Guyon2003} le atribuyen tres objetivos: mejorar el desempeño del predictor, obtener modelos más rápidos y económicos, y comprender mejor el proceso que generó los datos. \textcite{Kuhn2019} clasifican los métodos en filtros, que evalúan cada variable de forma independiente del modelo, y envolventes (\textit{wrappers}), que evalúan el modelo sobre subconjuntos candidatos de variables, ya sea mediante búsquedas voraces, como la eliminación recursiva, o mediante búsquedas globales; y advierten que, si la selección se realiza fuera del esquema de remuestreo, la estimación del error queda sesgada de forma optimista.

En el proyecto estas operaciones constituyen una etapa propia del pipeline de datos, que construye para cada proyecto un vector de variables a partir de sus atributos maestros y de agregados de sus registros transaccionales a una fecha de corte. Un requisito específico es que las mismas variables puedan calcularse tanto desde el histórico del portal anterior, que alimenta el entrenamiento, como desde el ERP nuevo, sobre el que se realiza la inferencia; el diseño de esa etapa se desarrolla en el Capítulo III\@.

% Candidatas @online para documento.bib (verificadas por curl el 2026-09-17; respaldan el soporte nativo de categóricas mencionado en 2.2.3):
% @online{XGBoostCategorical, title={Categorical Data}, organization={XGBoost developers}, url={https://xgboost.readthedocs.io/en/stable/tutorials/categorical.html}, urldate={2026-09-17}}  -- "Since version 1.5, XGBoost has support for categorical data" (parámetro enable_categorical).
% @online{SklearnEnsemble, title={Ensembles: Gradient boosting, random forests, bagging, voting, stacking -- Categorical Features Support}, organization={scikit-learn developers}, url={https://scikit-learn.org/stable/modules/ensemble.html}, urldate={2026-09-17}}  -- HistGradientBoostingClassifier/Regressor, parámetro categorical_features.
